\section{A rigidity certificate for uniform increment retyping}
\label{sec:increment-incidence}

Consecutive increment instructions commute, even when they increment the
same register.  This suggests changing which register is incremented by
each physical decoder row while keeping the net increment of every
consecutive increment block.  For the original Context246 layout, the
following finite certificate rules out every nontrivial retyping of this
specific kind.

The literal register program has 2,601 increment occurrences, partitioned
into 1,875 maximal consecutive increment blocks.  All 726 occurrences in
the interiors of these blocks have their immediately preceding increment
as their unique predecessor in the complete guarded instruction graph.
The initial control is not an interior occurrence.  Thus each block is
indeed a region in which all its increments execute consecutively after
entry; the count is not based only on a visual source listing.

Group increment occurrences by their existing physical source role
$(q,b)$, where $q$ is the first state to reach the parent decoder cell and
$b$ is its outgoing bit.  The 20 groups, in matrix-column order, are
\[
\begin{split}
 &(12,1),(26,0),(46,0),(50,0),(54,0),(55,0),(66,0),\\
 &(71,0),(81,0),(82,0),(90,0),(106,0),(114,0),(114,1),\\
 &(142,1),(147,0),(167,0),(170,0),(218,1),(233,1).
\end{split}
\]
Their original register labels in the same order are
\[
 (0,0,6,6,1,5,1,6,2,7,7,7,1,1,5,4,8,3,1,1).
\]
Each group has a unique original label across all its occurrences.

The independent verifier derives these roles from the literal transition
table.  For each parent path $p$, it runs from state 96 with the marker
and bits of $p$ on an otherwise zero tape, stopping on the first arrival
at cell $|p|$.  This permits the actual short left-and-right detours of
borrowed handlers; simply following rightward transition targets would
give an incorrect mapping.  In this data, all 2,338 distinct parent-path
simulations terminate within 40 literal transitions and stay between
cells 0 and 33.  A first-arrival simulation has not read any later cell,
so its returned state is independent of the following leaf bit, remaining
PC suffix, gap, and registers.  Its mapping agrees exactly with the
producer's 20 groups and every row of the incidence matrix.

Let $M\in\mathbb N^{1875\times20}$ be the incidence matrix: $M_{kg}$
counts occurrences of group $g$ in increment block $k$.  Direct integer
checks establish
\[
 \operatorname{rank}_{\mathbb Q}M=19,\qquad
 \ker_{\mathbb Q}M
 =\operatorname{span}_{\mathbb Q}
 \{e_{(114,0)}-e_{(114,1)}\}.
\]
The upper bound follows by checking that columns $(114,0)$ and $(114,1)$
are identical.  For the lower bound, the independently checked minor on
zero-based rows
\[
 0,1,2,3,4,5,6,10,12,13,14,16,17,26,33,162,210,1829,1834
\]
and zero-based columns, in this displayed order,
\[
 7,1,0,9,3,16,17,15,6,10,11,8,4,5,12,18,19,14,2
\]
has determinant 2.  These finite facts are sufficient for the stated
kernel; no floating-point rank decision is involved.

\begin{theorem}[Uniform-retyping rigidity]
Suppose a proposed change assigns a single register label to each of the
20 existing groups, retaining one increment per occurrence and retaining
the source topology and group membership.  If it preserves the total
register increment vector of every maximal increment block, every group
must retain its original register label.
\end{theorem}
\begin{proof}
Represent a label $r$ by the nonnegative unit vector $e_r\in\mathbb N^9$.
Let $X$ and $X_0$ be the $20\times9$ matrices whose rows are the proposed
and original group vectors.  The hypothesis is $MX=MX_0$.
For each register coordinate, its column of $X-X_0$ belongs to $\ker M$.
The kernel certificate therefore fixes all group vectors except possibly
those of $(114,0)$ and $(114,1)$, and fixes the sum of the latter two.
Both original labels are register 1, so that sum must remain $2e_1$.
The only two nonnegative unit vectors with sum $2e_1$ are $e_1$ and
$e_1$.  Thus these two labels are fixed as well.
\end{proof}

This is an exact negative result for a useful finite search dimension.
It is not a state lower bound.  It leaves open changes to decoder topology,
context-dependent row splitting, block boundaries, cross-block register
effects, and other simulation relations.  In particular, moving an
increment across a separate transfer loop lies outside its hypotheses.
The portable audit program is \texttt{check\_increment\_incidence.py};
it uses an independent source parser, a dense literal interpreter, and
exact rational elimination, and writes
\texttt{increment\_incidence\_audit.json} with the full nonzero minor.
